\documentclass[11pt]{article}
\usepackage[T1]{fontenc}
\usepackage[utf8]{inputenc}
\usepackage[a4paper,margin=1in]{geometry}
\usepackage{amsmath,amssymb,amsthm,mathtools}
\usepackage{microtype}
\usepackage{xcolor}
\usepackage{hyperref}
\usepackage{enumitem}
\usepackage{setspace}
\usepackage{lmodern}
\usepackage{booktabs}
\usepackage{array}
\hypersetup{colorlinks=true,linkcolor=blue!50!black,
urlcolor=blue!50!black,citecolor=blue!50!black}
\setlength{\parskip}{0.5em}
\setlength{\parindent}{0pt}
\onehalfspacing

\title{\vspace{-1.5em}\bfseries
From Local Curvature to the Weil Functional:\\[0.4em]
\large An Explicit Construction}
\author{Ulrich Tehrani}
\date{Research Note \textperiodcentered\ March 2026 \textperiodcentered\ v2}

\newtheorem*{lemma*}{Lemma}
\newtheorem*{proposition*}{Proposition}
\newtheorem*{corollary*}{Corollary}
\newtheorem*{remark*}{Remark}

\begin{document}
\maketitle

\begin{abstract}
Building on the curvature decomposition established in~\cite{paper1},
we construct an explicit family of admissible test functions
$g^*_{\sigma,\varepsilon} \in S_{ad}$ for the Weil explicit formula.
For this family we derive a prime-side identity of the form
\[
W(g^*_{\sigma,\varepsilon} * \check g^*_{\sigma,\varepsilon})
=
Z(g^*_{\sigma,\varepsilon}) - H_{\mathrm{local}}(\sigma,\kappa) + O(\varepsilon),
\]
where $Z(g^*_{\sigma,\varepsilon})$ denotes the corresponding zero sum.
We further introduce renormalized prime weights
$c_p^{\mathrm{ren}} > 0$, prove convergence of the associated series
\[
D = \sum_p \bigl(c_p^{\mathrm{ren}}\bigr)^2 \approx 9.471,
\]
and show that the renormalized prime-side scale is logarithmic in $\kappa$.
This isolates the remaining open problem as a single spectral inequality
relating $Z(g^*_{1/2,\varepsilon})$ to $H_{\mathrm{local}}(1/2,\kappa)$.
We also record the appearance of the Mertens constant as the natural scale
of the renormalized curvature remainder.
\end{abstract}

\section{Introduction and relation to Paper~1}

In~\cite{paper1} we studied the curvature field
\[
H_{\xi}(\sigma,t)
:= \frac{d^2}{d\sigma^2}\log|\xi(\sigma+it)|^2
= H_{\mathrm{local}}(\sigma,\kappa) + H_{\mathrm{dual}}(\sigma,t,\kappa),
\]
where $H_{\mathrm{dual}} := H_{\xi} - H_{\mathrm{local}}$ is the spectral
remainder encoding the zero contributions (sign convention as in~\cite{paper1},
equation~(5)).
where the truncated local curvature
\[
H_{\mathrm{local}}(\sigma,\kappa) = \psi_1(\sigma)
+ \sum_{p\le\kappa} V_p(\sigma)
\]
contains the archimedean term and the finite-place contributions. The prime-side pieces satisfy
$V_p(\sigma) \ge 0$, and on the critical line one has
$H_{\mathrm{local}}(1/2,\kappa) \asymp (\log\kappa)^2$.

The central open question of~\cite[\S7]{paper1} asks whether this local
positivity and critical divergence can be packaged into a global
Weil-type positivity statement. The present note addresses a concrete part of
that question by constructing an explicit family of admissible test functions
$g^*_{\sigma,\varepsilon} \in S_{ad}$ and identifying the corresponding
prime-side term in the Weil functional with the local curvature studied in
Paper~1.

The remaining open step is then isolated as a single spectral inequality;
see Section~\ref{sec:open}.

\section{Construction of the admissible test function}

\subsection{The raw family}

Recall from~\cite[Eq.~(3)]{paper1} that
\[
V_p(\sigma)
= 4(\log p)^2\,\frac{p^{-2\sigma}}{(1-p^{-2\sigma})^2}.
\]
Define
\[
c_{p,\sigma}
:= V_p(\sigma)^{1/2}
\cdot \frac{p^{1/4}}{(\log p)^{1/2}},
\]
and let
\[
\varphi_\varepsilon(x)
= \frac{1}{\varepsilon\sqrt{2\pi}}e^{-x^2/(2\varepsilon^2)}
\]
be the standard Gaussian bump. For fixed $\kappa$ set
\[
\widehat f_{\sigma,\varepsilon}(x)
= \sum_{p\le\kappa} c_{p,\sigma}\,\varphi_\varepsilon(x-\log p).
\]
The admissible test function is obtained by antisymmetrization:
\begin{equation}\label{eq:gstar}
\widehat g^*_{\sigma,\varepsilon}(x)
:= \widehat f_{\sigma,\varepsilon}(x)
- \widehat f_{\sigma,\varepsilon}(-x).
\end{equation}

\begin{lemma*}[$g^*_{\sigma,\varepsilon} \in S_{ad}$]
For every $\sigma>0$ and $\varepsilon>0$, the function
$g^*_{\sigma,\varepsilon}$ belongs to the admissible class $S_{ad}$ used in
Lagarias' formulation of Weil positivity.
\end{lemma*}

\textit{Proof.}
The Gaussian bumps are Schwartz, and antisymmetrization preserves the
Schwartz class. Hence $g^*_{\sigma,\varepsilon}$ and
$\widehat g^*_{\sigma,\varepsilon}$ are rapidly decreasing. Moreover,
$\widehat g^*_{\sigma,\varepsilon}(0)=0$ by antisymmetry, and likewise
$g^*_{\sigma,\varepsilon}(0)=0$. These are the required vanishing conditions
at the origin. Finally, the Gaussian decay ensures absolute convergence of the
prime-side and zero-side expressions appearing in the Weil functional for this
family. \hfill$\square$

\section{Prime-side identity for the Weil functional}

\begin{proposition*}[Prime-side identity]
For the family $g^*_{\sigma,\varepsilon}$ constructed above, one has
\begin{equation}\label{eq:weil-decomp}
W(g^*_{\sigma,\varepsilon} * \check g^*_{\sigma,\varepsilon})
=
Z(g^*_{\sigma,\varepsilon}) - H_{\mathrm{local}}(\sigma,\kappa) + O(\varepsilon),
\end{equation}
where
$Z(g^*_{\sigma,\varepsilon}) = \sum_{\rho}
|\widehat g^*_{\sigma,\varepsilon}(\rho)|^2$
runs over the nontrivial zeros of $\zeta(s)$.
\end{proposition*}

\textit{Proof sketch.}
Apply the Weil explicit formula in the Lagarias framework. The constant
term vanishes because $\widehat g^*_{\sigma,\varepsilon}(0)=0$.
By construction of the coefficients $c_{p,\sigma}$, the prime-side quadratic
form converges to $H_{\mathrm{local}}(\sigma,\kappa)$ as $\varepsilon\to 0$.
The error term $O(\varepsilon)$ arises from the Gaussian smoothing of the
prime-power contributions and is not made uniform in $\kappa$ in the present
note. \hfill$\square$

\begin{remark*}
Equation~\eqref{eq:weil-decomp} is the key bridge between the curvature field
of Paper~1 and the Weil positivity framework. The archimedean remainder is
exactly zero since $\widehat g^*_{\sigma,\varepsilon}(0) = 0$ (condition S2 of
$S_{ad}$).
\end{remark*}

\section{Renormalized weights and convergence}

\subsection{Positivity of the renormalized coefficients}

At $\sigma=1/2$ define
\[
f_p := V_p(1/2) - \frac{4(\log p)^2}{p}.
\]
A direct calculation gives
\begin{equation}\label{eq:fp}
f_p
= \frac{4(\log p)^2(2p-1)}{p(p-1)^2} > 0
\qquad (p\ge 2).
\end{equation}

\begin{lemma*}[Renormalized coefficients are positive]
For every prime $p\ge 2$, the quantity $f_p$ in~\eqref{eq:fp}
is strictly positive.
\end{lemma*}

\textit{Proof.}
The numerator $2p-1>0$ and the denominator $p(p-1)^2>0$
for every prime $p\ge 2$. \hfill$\square$

The renormalized coefficients are
\begin{equation}\label{eq:cren}
c_p^{\mathrm{ren}}
:= f_p^{1/2}\,\frac{p^{1/4}}{(\log p)^{1/2}} > 0.
\end{equation}
Numerical values for small primes:

\begin{center}
\begin{tabular}{ccccc}
\toprule
$p$ & $V_p(1/2)$ & $4(\log p)^2/p$ & $f_p$
    & $[c_p^{\mathrm{ren}}]^2$ \\
\midrule
2  & 5.545 & 0.962 & 4.583 & 2.883 \\
3  & 3.625 & 1.600 & 2.025 & 2.012 \\
5  & 3.213 & 2.058 & 1.155 & 1.166 \\
7  & 3.019 & 2.301 & 0.718 & 0.781 \\
11 & 2.830 & 2.570 & 0.260 & 0.439 \\
\bottomrule
\end{tabular}
\end{center}

\subsection{Convergence of the renormalized energy}

\begin{lemma*}[Convergence]
The series
\begin{equation}\label{eq:D}
D := \sum_p \bigl[c_p^{\mathrm{ren}}\bigr]^2
   = \sum_p \frac{4(\log p)^2(2p-1)}{p(p-1)^2}
\end{equation}
converges absolutely. Numerically $D \approx 9.471$
(primes below $10^4$; tail $<0.009$).
\end{lemma*}

\textit{Proof.}
For large $p$, $[c_p^{\mathrm{ren}}]^2 \sim 8(\log p)^2/p^2$.
Since $\sum_n (\log n)^2/n^2 <\infty$, convergence follows
by comparison. \hfill$\square$

\section{The renormalized scale and the Mertens constant}

By the prime number theorem,
\[
H_{\mathrm{local}}(1/2,\kappa)
= 2(\log\kappa)^2 + A\gamma_M\log\kappa + C_0 + o(\log\kappa),
\]
where $\gamma_M \approx 0.2615$ is the Mertens constant,
\[
\gamma_M = \gamma + \sum_p\!\Bigl[\log\tfrac{p}{p-1}-\tfrac1p\Bigr],
\]
and $\gamma \approx 0.5772$ is the Euler--Mascheroni constant.
The renormalized remainder
\[
H_{\mathrm{local}}^{\mathrm{ren}}(1/2,\kappa)
:= H_{\mathrm{local}}(1/2,\kappa) - 2(\log\kappa)^2
\sim A\gamma_M\log\kappa
\]
lives on the logarithmic scale, and the Mertens constant
is its natural coefficient.

\begin{remark*}
The Euler--Mascheroni constant $\gamma$ appears simultaneously
in $\gamma_M$ (prime side, via the curvature remainder) and in the
first Li coefficient
$\lambda_1 = 1 + \gamma/2 - \log 2 - (\log\pi)/2 \approx 0.0231$
(zero side). This structural coincidence reflects the universal role
of $\gamma$ as the renormalization constant governing the transition
from discrete prime sums to continuous logarithmic integrals.
\end{remark*}

\begin{remark*}
The renormalized construction clarifies the scale of the prime-side
contribution, but does not by itself resolve the spectral inequality
\eqref{eq:openineq}.
\end{remark*}

\section{The remaining spectral inequality}\label{sec:open}

For the family $g^*_{1/2,\varepsilon}$, Weil positivity reduces to
\begin{equation}\label{eq:openineq}
Z(g^*_{1/2,\varepsilon})
\;\ge\;
H_{\mathrm{local}}(1/2,\kappa).
\end{equation}

\textbf{Open question.}
Is there an unconditional lower bound
\[
Z(g^*_{1/2,\kappa})
= \sum_{\rho}\bigl|\widehat g^*_{1/2,\kappa}(\rho)\bigr|^2
\;\ge\;
C\cdot(\log\kappa)^2
\]
for some constant $C>0$, where the sum runs over nontrivial zeros of~$\zeta$?

A Selberg-type heuristic suggests $Z(g^*)\sim D\log\kappa\log\log\kappa$
with $D\approx 9.471$. An unconditional proof of even
$Z\ge C(\log\kappa)^2$ is presently unavailable.

\begin{remark*}
By Lagarias~\cite{lagarias},
$\mathrm{RH} \iff W(f*\check f)\ge 0$ for all $f\in S_{ad}$.
A positive answer to the open question above would establish
$W(g^*_{1/2,\kappa}*\check g^*_{1/2,\kappa})\ge 0$ for the specific
family constructed here, which is a necessary but not sufficient
condition for RH in the absence of a density argument for this
family in~$S_{ad}$.
\end{remark*}

\section*{Computational note}

All numerical claims are reproducible. Code is available at~\cite{github}.
The Gaussian parameter $\varepsilon = 0.05$ is required for meaningful
numerical experiments with the zero sum (for $\varepsilon = 0.5$ one has
$e^{-\varepsilon^2\gamma_1^2} \approx 10^{-22}$, effectively
suppressing all zero contributions).

\begin{thebibliography}{99}

\bibitem{paper1}
U.~Tehrani,
\emph{A Curvature Decomposition of the Explicit Formula
for the Riemann Zeta Function},
Research Note, March 2026, v6,
\href{https://doi.org/10.5281/zenodo.19025598}{%
doi:10.5281/zenodo.19025598}.

\bibitem{github}
U.~Tehrani,
\emph{AnalysisLab: Curvature and the Weil Functional},
GitHub repository, 2026,
\url{https://github.com/utehrani/analysislab-nt}.

\bibitem{lagarias}
J.~C.~Lagarias,
\emph{Li coefficients for automorphic L-functions},
Ann.\ Inst.\ Fourier \textbf{57}(5) (2007), 1689--1740.

\bibitem{li}
X.-J.~Li,
\emph{The positivity of a sequence of numbers
and the Riemann hypothesis},
J.\ Number Theory \textbf{65} (1997), 325--333.

\bibitem{weil}
A.~Weil,
\emph{Sur les ``formules explicites'' de la
th\'eorie des nombres premiers},
Comm.\ S\'em.\ Math.\ Univ.\ Lund (1952), 252--265.

\bibitem{mertens}
F.~Mertens,
\emph{Ein Beitrag zur analytischen Zahlentheorie},
J.\ Reine Angew.\ Math.\ \textbf{78} (1874), 46--62.

\bibitem{selberg}
A.~Selberg,
\emph{Contributions to the theory of the Riemann zeta-function},
Arch.\ Math.\ Naturvid.\ \textbf{48}(5) (1946), 89--155.

\bibitem{connes}
A.~Connes and C.~Consani,
\emph{Weil positivity and trace formula, the archimedean place},
Sel.\ Math.\ (N.S.) \textbf{27} (2021), Article~77.

\bibitem{edwards}
H.~M.~Edwards,
\emph{Riemann's Zeta Function},
Dover, 1974.

\end{thebibliography}

\section*{Acknowledgment}

The author thanks Jean-Fran\c{c}ois Burnol for confirming that the
curvature decomposition of~\cite{paper1} does not appear in the
literature known to him.

\end{document}
